\section{Mediciones y estimación}\label{sec:mediciones}

El modelo de la Sección~\ref{sec:modelo-general} solo es útil si sus
estados y parámetros pueden recuperarse de las mediciones disponibles.
Esta sección define el mapa de observación del sistema instrumentado
(Sección~\ref{sec:diseno}), analiza qué es observable e identificable
con esa instrumentación, y establece el protocolo de caracterización
previa y la estrategia de estimación que alimenta el gemelo digital.

\subsection{Vector de mediciones y mapa de observación}\label{sec:mapa-obs}

La instrumentación del diseño aporta tres grupos de señales, muestreadas
a $0{,}1$~Hz ($10$~s entre muestras, Observación~\ref{rem:muestreo-cierres})
y registradas con marca temporal común:

\begin{description}
  \item[Gases y aire interno (outlet).] Temperatura $T$ y humedad relativa
  $H$ (SHT35), CO$_2$ (NDIR), CH$_4$ y O$_2$ en la corriente de salida:
  \[
    \mathbf{y}_{\mathrm{gas}}(t) =
    \bigl(T,\; H,\; c_{\mathrm{CO_2}},\; c_{\mathrm{CH_4}},\;
    c_{\mathrm{O_2}}\bigr)(t).
  \]
  Bajo el supuesto de mezcla homogénea (A1), estas lecturas son los
  propios estados del aire interno: el mapa de observación es aquí la
  identidad, más ruido de medición.
  \item[Lecho.] El sensor capacitivo insertado en el lecho aporta
  $(\theta_s^{\mathrm{sen}}, T_s^{\mathrm{sen}})$. La temperatura lee el
  estado $T_s$ directamente; el canal capacitivo requiere calibración
  gravimétrica para relacionarse con $\theta_s$
  (Sección~\ref{sec:caracterizacion}):
  \[
    \theta_s^{\mathrm{sen}}(t) = h_\theta\bigl(\theta_s(t)\bigr) + \varepsilon_\theta,
    \qquad
    T_s^{\mathrm{sen}}(t) = T_s(t) + \varepsilon_T.
  \]
  \item[Masa del lote.] La celda de carga registra $m_{\mathrm{LC}}(t)$,
  la masa total del lote, sujeta a la restricción de consistencia global
  de la Ec.~\eqref{eq:celda-carga}: su tasa de decrecimiento debe
  igualar el flujo neto de masa por el \emph{outlet}.
\end{description}

Además, las variables de supervisión (H$_2$S, NH$_3$) se registran con
fines de alarma sin entrar al mapa de observación del modelo
(Sección~\ref{sec:sistema}), y las condiciones del \emph{inlet}
$(T_{\mathrm{in}}, w_{\mathrm{in}}, c_{\mathrm{O_2,in}})$ se miden como
forzamientos conocidos, no como estados.

\subsection{Observabilidad con $(T, H, \mathrm{CO_2}, \mathrm{CH_4}, \mathrm{O_2})$}\label{sec:observabilidad}

El vector de estado tiene dimensión $12$
(Ec.~\eqref{eq:estado-completo}). Su observabilidad se analiza por
grupos de compartimentos:

\begin{description}
  \item[Aire interno $(T, w, c_{\mathrm{CO_2}}, c_{\mathrm{CH_4}}, c_{\mathrm{O_2}})$.]
  Observable directamente: cinco estados medidos por identidad, con $w$
  obtenido de $H$ y $T$ vía psicrometría.
  \item[Lecho $(\theta_s, T_s)$.] Observable con el sensor capacitivo
  calibrado y el balance de agua~\eqref{eq:agua-lecho} como verificación
  cruzada.
  \item[Sustrato $S$.] No medible en línea de forma directa; se estima
  por el balance~\eqref{eq:sustrato-general} una vez conocidas las tasas,
  con ancla en la celda de carga y en la determinación gravimétrica
  inicial.
  \item[Larva promedio $(A, B, L)$ y población $N$.] Estados latentes:
  ningún sensor los lee. Su única vía de observación son las tasas
  agregadas $R_i$ medidas en los cierres
  (Observación~\ref{rem:cierres}), que por~\eqref{eq:rco2-general}
  combinan $mB$, $Y_B r_B$ y $Y_L r_L$. El par $(\mathrm{CO_2}, \mathrm{O_2})$
  aporta además el cociente respiratorio medido~\eqref{eq:rq-medido},
  que discrimina el sustrato metabólico oxidado.
\end{description}

En síntesis: siete estados son medibles de forma directa o semidirecta,
y cinco $(A, B, L, N, S)$ son estimables solo a través del modelo. Esta
estructura justifica la estrategia de estimación de la
Sección~\ref{sec:estimacion}: los estados medidos anclan, y las tasas de
los cierres informan a los latentes.

\subsection{Identificabilidad: separación larva--microbiota}\label{sec:identificabilidad}

Las señales del \emph{outlet} son superposiciones
(Ec.~\eqref{eq:co2-superpuesto}): sin información adicional, la
contribución larval y la microbiana no son separables, y el
$RQ_{\mathrm{medido}}$ no es interpretable fisiológicamente
(Ec.~\eqref{eq:rq-medido}). El protocolo resuelve esta identificabilidad
con tres mecanismos complementarios:

\begin{enumerate}
  \item \textbf{Línea base microbiana temprana.} El primer cierre de
  ventilación, ejecutado dentro de las primeras $24$~h tras la carga,
  mide una señal dominada por $r^{\mathrm{mic}}$ (larvas neonatas de
  masa despreciable). Es la estimación inicial de la actividad
  microbiana del sustrato.
  \item \textbf{Lote testigo sin larvas.} Una corrida completa con el
  mismo sustrato, humedad inicial y calendario de cierres, pero
  $N_0 = 0$, mide $r^{\mathrm{mic}}_{\mathrm{CO_2}}$ y
  $r^{\mathrm{mic}}_{\mathrm{O_2}}$ de forma independiente en todo el
  horizonte. Este tratamiento es parte obligatoria del protocolo
  experimental: sin él, la descomposición
  de~\eqref{eq:co2-superpuesto} queda indeterminada.
  \item \textbf{Cierre post-cosecha.} Un cierre adicional tras retirar
  las prepupas mide la respiración residual del frass, que acota la
  actividad microbiana al final del ciclo.
\end{enumerate}

\begin{remark}[El CH$_4$ como testigo de anaerobiosis]
La metanogénesis es casi exclusivamente microbiana
(Observación~\ref{rem:microbiano}), de modo que el CH$_4$ no sufre el
problema de superposición: su señal informa directamente la fracción
anaerobia latente $\xi$ (Observación~\ref{rem:xi}) y valida la firma de
terminación propuesta en la Observación~\ref{rem:tf}.
\end{remark}

\subsection{Caracterización previa del sistema}\label{sec:caracterizacion}

Antes de cargar el primer lote, el sistema vacío y los sensores se
caracterizan en cuatro pruebas, cuyos resultados son parámetros del
modelo o de la cadena de medición:

\begin{enumerate}
  \item \textbf{Prueba de hermeticidad} (Observación~\ref{rem:fugas}):
  estimación de $k_{\mathrm{fuga}}$ por decaimiento de CO$_2$ inyectado.
  Si $k_{\mathrm{fuga}}$ no es despreciable, se incorpora como término
  correctivo en los balances de gases.
  \item \textbf{Volumen de aire efectivo.} $V_{\mathrm{aire}}$ se mide
  por dilución: inyección de un pulso conocido de CO$_2$ en la cámara
  cerrada con el lecho simulado, y lectura del incremento de
  concentración en estado estacionario. Este valor, no el geométrico,
  es el que entra en~\eqref{eq:modo-cerrado}.
  \item \textbf{Calibración de sensores.} El NDIR de CO$_2$ se calibra
  en dos puntos (aire ambiente $\sim 420$~ppm y gas patrón); el sensor
  capacitivo del lecho se calibra contra determinaciones gravimétricas
  de $\theta_s$ en el rango de operación.
  \item \textbf{Validación contra cromatografía.} En cada corrida se
  toman muestras por el septum en instantes seleccionados (inicio de
  cierre, mitad y fin) para análisis por cromatografía de gases. La
  comparación GC--sensor en línea cuantifica el sesgo y la deriva de
  los sensores de CO$_2$ y CH$_4$, y documenta la calidad metrológica
  del experimento.
\end{enumerate}

Las mismas pruebas se repiten en la recalibración entre réplicas
secuenciales (Observación~\ref{rem:escalado}).

\subsection{Estimación de estados y parámetros}\label{sec:estimacion}

La estimación se organiza en dos niveles, coherentes con la distinción
entre datos de identificación y datos de validación
(Observación~\ref{rem:id-lazo}):

\paragraph{Nivel 1: tasas instantáneas.} En cada cierre
$[t_k^c, t_k^c + \Delta]$, las pendientes de acumulación se ajustan por
mínimos cuadrados sobre los $\sim 9$ puntos disponibles
(Observación~\ref{rem:muestreo-cierres}) y entregan las tasas agregadas
\[
  R_i(t_k) = \sigma_i\, V_{\mathrm{aire}}\,
  \widehat{\left.\dot c_i\right|_{[t_k^c,\, t_k^c + \Delta]}},
  \qquad i \in \{\mathrm{CO_2}, \mathrm{O_2}, \mathrm{CH_4}\},
\]
independientes de $Q$ y sin modelado fisiológico. La ventana $\Delta$ se
ajusta día a día según las cotas~\eqref{eq:delta-min}
y~\eqref{eq:delta-max}, evaluadas con las tasas del día anterior: es la
primera función del gemelo digital.

\paragraph{Nivel 2: estados y parámetros del modelo.} Las series
$R_i(t_k)$, las señales continuas de la Sección~\ref{sec:mapa-obs} y las
condiciones iniciales conocidas alimentan el ajuste de los parámetros
$\boldsymbol{\theta}$ (fisiológicos, microbianos, térmicos e hídricos)
y la reconstrucción de los estados latentes $(A, B, L, N, S)$. La
vía prevista es una red neuronal informada por física (PINN), con las
ecuaciones de la Sección~\ref{sec:modelo-general} como restricciones de
la pérdida; un filtro de Kalman extendido o de ensamble sobre la forma
unificada ventilado--cerrado sirve como referencia comparativa. Los
datos de operación regulada normal se reservan para validación: el
modelo ajustado debe predecir las concentraciones entre cierres sin
haberlas usado en el ajuste.

\paragraph{Criterio de éxito.} La estimación se considera exitosa si
(i) los residuos de validación de $c_{\mathrm{CO_2}}$ y
$c_{\mathrm{O_2}}$ son compatibles con el ruido de los sensores,
(ii) la restricción de la celda de carga~\eqref{eq:celda-carga} se
cumple dentro de su tolerancia, y (iii) los parámetros fisiológicos
estimados caen en los rangos reportados para el modelo de referencia
\citep{eriksenDynamicModellingFeed2022,eriksenMetabolicPerformanceFeed2024}.

\subsection{Dashboard del sistema monitoreado}\label{sec:dashboard}

El tablero web es la interfaz del gemelo digital y despliega, como
mínimo, cuatro componentes:

\begin{enumerate}
  \item \textbf{Series en tiempo real} de $\mathbf{y}_{\mathrm{gas}}$,
  $(\theta_s^{\mathrm{sen}}, T_s^{\mathrm{sen}})$ y $m_{\mathrm{LC}}$,
  con los eventos de cierre marcados sobre las series de gases.
  \item \textbf{Tasas estimadas} $R_i(t_k)$ del Nivel 1, superpuestas
  con la predicción del modelo ajustado (Nivel 2): la brecha entre
  ambas es el diagnóstico en línea del gemelo.
  \item \textbf{Indicadores fisiológicos derivados:}
  $RQ_{\mathrm{medido}}$ y la razón CH$_4$/CO$_2$, con la firma de
  terminación de la Observación~\ref{rem:tf} evaluada en tiempo real.
  \item \textbf{Supervisión y alarmas}: H$_2$S y NH$_3$ contra umbrales,
  estado de los actuadores $(Q, P_{\mathrm{pel}}, P_{\mathrm{hum}})$ y
  cumplimiento de las cotas de seguridad de O$_2$ durante los cierres.
\end{enumerate}

La persistencia es una base de datos en la nube alimentada por los
ESP32 (Sección~\ref{sec:diseno}); cada réplica queda archivada como un
conjunto completo de series, eventos y metadatos del lote (cepa, edad y
cantidad de neonatos, $\theta_{s,0}$, $S_0$), reutilizable en el
análisis estadístico posterior.
